\documentclass[11pt]{article}
\usepackage[margin=1in]{geometry}
\usepackage{amsmath,amssymb,amsthm}
\usepackage[hidelinks]{hyperref}

\newtheorem{theorem}{Theorem}
\newtheorem{lemma}{Lemma}
\newtheorem{corollary}{Corollary}

\title{A Polynomial-Space $O(1.4289^n)$ Algorithm for\\ 3-Path Vertex Cover on Subcubic Graphs}
\author{}
\date{}

\begin{document}
\maketitle

\begin{abstract}
We observe that a greedy induced matching yields a dissociation set of size at least
$n/3$ in every subcubic graph. Hence the minimum 3-path vertex cover has size at most
$2n/3$. Combining this bound with the polynomial-space parameterized algorithm of
\v{C}erven\'y and Such\'y (arXiv:2111.05896), which runs in $O^*(1.708^k)$ time, gives
an $O^*(1.708^{2n/3}) = O(1.4289^n)$-time polynomial-space algorithm for 3-path vertex
cover on subcubic graphs. This improves the previous polynomial-space bound
$O(1.4613^n)$ of Chang et al.\ (2018). It does not improve the exponential-space bound
$O(1.3659^n)$.
\end{abstract}

\section{Definitions and known results}

A \emph{3-path} is a path on three vertices. A \emph{3-path vertex cover} is a set
$C \subseteq V$ such that $G-C$ contains no 3-path, equivalently $G-C$ has maximum
degree at most $1$. A \emph{dissociation set} is a set $S \subseteq V$ with
$\Delta(G[S]) \le 1$. The two notions are complementary: $S$ is a dissociation set if
and only if $V \setminus S$ is a 3-path vertex cover. Let $k^*(G)$ be the minimum size
of a 3-path vertex cover of $G$.

Known exact algorithms for an $n$-vertex graph are
$O^*(1.5171^n)$ (Kardoš, Katrenič, and Schiermeyer),
$O^*(1.4658^n)$ (Chang et al.),
$O^*(1.4613^n)$ time and polynomial space (Chang et al., 2018),
$O^*(1.3659^n)$ time with exponential space (Xiao and Kou, 2017).

For the parameterized version, the running time was improved successively to
$O^*(2^k)$ (Tu), $O^*(1.8127^k)$ (Katrenič), $O^*(1.7485^k)$ (Xiao and Kou, polynomial
space), $O^*(1.713^k)$ (Tsur), and $O^*(1.708^k)$ (\v{C}erven\'y and Such\'y, 2021,
branch-and-reduce, polynomial space).

\section{A dissociation set of size at least $n/3$}

\begin{lemma}
Every subcubic graph $G$ on $n$ vertices has a dissociation set of size at least
$n/3$.
\end{lemma}

\begin{proof}
Construct a set $S$ greedily. Start with $S = \emptyset$. While $G$ has an edge, choose
an edge $uv$, add $u$ and $v$ to $S$, and delete $N[u] \cup N[v]$ from $G$. When no
edge remains, add all remaining (isolated) vertices to $S$.

We claim $G[S]$ has maximum degree at most $1$. Every time an edge $uv$ is chosen, all
neighbours of $u$ and $v$ are deleted; hence no vertex of $S$ outside $\{u,v\}$ is
adjacent to $u$ or $v$, and no edge between the endpoints of two different chosen
edges survives. The only edges inside $S$ are the chosen matching edges; the vertices
added at the end are isolated. Thus $G[S]$ is a matching plus isolated vertices.

For the size bound, let $t$ be the number of chosen edges. Each step deletes the two
endpoints (which are in $S$) and at most $2(\Delta-1) \le 4$ other neighbours, so the
number of deleted vertices that are not in $S$ is at most $4t$. Hence
\[
  n \;\le\; |S| + 4t .
\]
Since the $2t$ endpoints are distinct vertices of $S$, we have $2t \le |S|$, and
therefore
\[
  n \;\le\; |S| + 2|S| \;=\; 3|S| ,
\]
so $|S| \ge n/3$.
\end{proof}

\begin{corollary}
Every subcubic graph satisfies $k^*(G) \le 2n/3$.
\end{corollary}

\begin{proof}
By Lemma 1 there is a dissociation set $S$ of size at least $n/3$. Its complement
$V \setminus S$ is a 3-path vertex cover of size at most $2n/3$.
\end{proof}

\section{The algorithm}

\begin{enumerate}
\item Compute the greedy set $S$ of Lemma 1 and let $k_0 = n - |S|$. By Corollary 2,
      $k^*(G) \le k_0 \le 2n/3$.
\item For $k = 0,1,\dots,k_0$, run the polynomial-space decision algorithm
      \texttt{P3VC}$(G,k)$ of \v{C}erven\'y and Such\'y.
\item Return the smallest $k$ for which \texttt{P3VC}$(G,k)$ answers yes.
\end{enumerate}

\textbf{Correctness.} \texttt{P3VC} is exact, and $k^*(G) \le k_0$; the loop therefore
returns $k^*(G)$.

\textbf{Time.} The $k$-th call takes $O^*(1.708^k)$ time. The sum over
$k=0,\dots,k_0$ is dominated by the last term, hence
\[
  O^*(1.708^{k_0}) \;\le\; O^*(1.708^{2n/3}) \;=\; O(1.4289^n).
\]
($1.708^{2/3} = 1.42886\ldots$.)

\textbf{Space.} The algorithm of \v{C}erven\'y and Such\'y is a branch-and-reduce
algorithm and uses polynomial space; the loop only stores the current instance. Hence
the whole algorithm uses polynomial space.

\section{Comparison}

\begin{center}
\begin{tabular}{lll}
\hline
Algorithm & Time & Space \\
\hline
Kardoš et al.\ 2011 & $O^*(1.5171^n)$ & polynomial \\
Chang et al. & $O^*(1.4658^n)$ & polynomial \\
Chang et al.\ 2018 & $O^*(1.4613^n)$ & polynomial \\
Xiao--Kou 2017 & $O^*(1.3659^n)$ & exponential \\
\textbf{This note} & $\boldsymbol{O(1.4289^n)}$ & \textbf{polynomial} \\
\hline
\end{tabular}
\end{center}

The new bound improves the polynomial-space bound $O(1.4613^n)$ to $O(1.4289^n)$.
It does not improve the exponential-space bound $O(1.3659^n)$.

\section{Remarks and limitations}

\begin{itemize}
\item The result is specific to graphs of maximum degree 3.
\item The improvement over the previous polynomial-space bound is about $2.2\%$.
\item The argument is black-box with respect to the parameterized algorithm. If the
      parameterized bound $1.708^k$ is improved, the same construction improves
      accordingly: a parameterized bound $c^k$ and the bound $k \le 2n/3$ give
      $c^{2n/3}$.
\item A larger constructive dissociation set would improve the bound further. The greedy
      construction gives $n/3$. A bound of $2n/5$ (i.e.\ $k \le 3n/5$) would give
      $1.708^{3/5} = 1.3788^n$. Reaching the exponential-space bound $1.3659^n$ would
      require $k/n \le 0.5825$, i.e.\ a constructive dissociation set of size at least
      $0.4175n$.
\end{itemize}

\section{Novelty check}

A web search for the bound $1.4289$ and related queries did not find a prior
$1.4289^n$ bound. The relevant prior work is the parameterized algorithm of
\v{C}erven\'y and Such\'y ($O^*(1.708^k)$, branch-and-reduce), the exponential-space
bound $O^*(1.3659^n)$ of Xiao and Kou, and the polynomial-space bound
$O^*(1.4613^n)$ of Chang et al.\ (2018). The hybrid used here is elementary and may be
folklore; novelty should be verified against the latest literature before submission.

\begin{thebibliography}{9}

\bibitem{cs21}
R.~\v{C}erven\'y and O.~Such\'y.
Generating faster algorithms for $d$-Path Vertex Cover.
arXiv:2111.05896, 2021.

\bibitem{tsur}
D.~Tsur.
Parameterized algorithm for 3-path vertex cover.
arXiv:1809.02636, 2018.

\bibitem{chang18}
M.-S.~Chang, L.-H.~Chen, L.-J.~Hung, Y.-Z.~Liu, P.~Rossmanith, and S.~Sikdar.
Moderately exponential time algorithms for the maximum bounded-degree-1 set problem.
\emph{Discrete Applied Mathematics}, 2018.

\bibitem{xk17}
M.~Xiao and S.~Kou.
Exact algorithms for the maximum dissociation set and minimum 3-path vertex cover
problems.
\emph{Theoretical Computer Science} 657 (2017), 86--97.

\bibitem{xk16}
M.~Xiao and S.~Kou.
Kernelization and Parameterized Algorithms for 3-Path Vertex Cover.
arXiv:1608.07022, 2016.

\bibitem{kks11}
F.~Kardoš, J.~Katrenič, and I.~Schiermeyer.
On computing the minimum 3-path vertex cover and dissociation number of graphs.
\emph{Theoretical Computer Science} 412 (2011), 7009--7017.

\end{thebibliography}

\end{document}
